% !TEX root = ../../../main.tex
% Sources: LEC03, IMG_9131--IMG_9133; e argument from LEC01, IMG_9111--IMG_9114.
% Geometric estimates and the positional-representation proof bridge prerequisites.
\section{Denseness, closure, and completion}
\label{sec:analysis-i:denseness}

\subsection{Upper and lower limits}
Before returning to sets, recall that $x_n\to L$ means that all sufficiently
late terms lie within every prescribed distance of $L$.
\begin{definition}
  A \emph{subsequence} of $(x_n)$ is a sequence $(x_{n_k})$ with
  $n_1<n_2<\cdots$. A real number $\ell$ is a \emph{limit point of the
  sequence} if some subsequence converges to $\ell$.
\end{definition}
We form a subsequence by keeping infinitely many terms in their original
order. The index $k$ counts the selected terms, while $n_k$ records their
positions in the original sequence. Since $n_k\geq k$, those positions
go to infinity as $k$ does. Thus a subsequence samples arbitrarily late
terms, rather than repeatedly choosing an early term.

If $x_n\to L$, every subsequence also converges to $L$. Indeed, once
all terms with $n\geq N$ lie within $\varepsilon$ of $L$, the selected
terms with $k\geq N$ do too, because $n_k\geq k\geq N$.
A convergent subsequence describes the behaviour of its selected terms.
Establishing a limit for the whole sequence requires controlling every
sufficiently late term.

A sequence may have several limit points even though the whole sequence
does not converge. For a bounded sequence, two useful quantities summarize
what remains after its first $n-1$ terms are discarded,
\[
  \alpha_n=\inf\{x_k\mid k\geq n\},\qquad
  \beta_n=\sup\{x_k\mid k\geq n\}.
\]
The sequence $(\alpha_n)$ is nondecreasing, and $(\beta_n)$ is nonincreasing.
Removing a term from a tail cannot lower its infimum or raise its supremum.
Both sequences are bounded, so their limits exist. Write
\[
  \ell_-=\liminf_{n\to\infty}x_n=\lim_{n\to\infty}\alpha_n,
  \qquad
  \ell_+=\limsup_{n\to\infty}x_n=\lim_{n\to\infty}\beta_n.
\]
These are also written $\underline{\lim}\,x_n$ and $\overline{\lim}\,x_n$.
The intervals $I_n=[\alpha_n,\beta_n]$ are nested, and $I_n$ contains
every term $x_k$ with $k\geq n$. Their endpoints approach $\ell_-$ and
$\ell_+$. Discarding finitely many initial terms changes neither limit.

\autoref{fig:tail-bounds-and-limit-points} shows this for
$x_n=(-1)^n(1+1/n)$. Choosing $n_k=2k-1$ keeps the odd terms, while
choosing $n_k=2k$ keeps the even terms. Taking $k\to\infty$ gives
\[
  x_{2k-1}=-1-\frac{1}{2k-1}\longrightarrow-1,\qquad
  x_{2k}=1+\frac{1}{2k}\longrightarrow1.
\]
The original sequence cannot converge, since its two subsequences would
then have to share the same limit. These are the only limit points,
since $|x_n|\to1$ forces
every subsequence limit $\ell$ to satisfy $|\ell|=1$.
The tail intervals shrink towards $[-1,1]$, but the numbers strictly
between $-1$ and $1$ are not limit points of this sequence.
\begin{figure*}[htbp]
  \centering
  % Exact sequence terms and tail intervals, not a schematic rescaling.
\begin{tikzpicture}[analysis diagram]
  \begin{scope}[x=4mm,y=8mm]
    \node at (6.5,2.7) {\textit{(a) Sequence terms}};
    \node at (6.5,2.2) {$x_n=(-1)^n(1+1/n)$};
    \draw[analysis guide] (0,1) -- (12.8,1);
    \draw[analysis guide] (0,-1) -- (12.8,-1);
    \draw[->] (0,-2.3) -- (0,1.9) node[above=5pt] {$x_n$};
    \draw[->] (0,0) -- (13.4,0) node[right=5pt] {$n$};
    \node[left=5pt] at (0,1) {$1$};
    \node[left=5pt] at (0,-1) {$-1$};
    \node[below left=5pt] at (0,0) {$0$};
    \foreach \n in {4,8,12} {
      \draw (\n,0.07) -- (\n,-0.07);
      \node[below=5pt] at (\n,-0.07) {$\n$};
    }
    \foreach \n in {1,...,12} {
      \fill (\n,{(-1)^\n*(1+1/\n)}) circle[radius=1.6pt];
    }
  \end{scope}

  \begin{scope}[xshift=10.7cm,x=12mm,y=8mm]
    \node at (-0.25,2.7) {\textit{(b) Bounds on successive tails}};
    \foreach \n/\lower/\upper/\row/\leftlabel/\rightlabel in {
      1/-2/1.5/1.4/{-2}/{3/2},
      3/{-4/3}/1.25/0.35/{-4/3}/{5/4},
      9/{-10/9}/1.1/-0.7/{-10/9}/{11/10}} {
      \node[anchor=east,interval label] at (-2.6,\row) {$I_{\n}$};
      \draw (\lower,\row) -- (\upper,\row);
      \node[interval closed] at (\lower,\row) {};
      \node[interval closed] at (\upper,\row) {};
      \node[above=5pt,interval label] at (\lower,\row) {$\leftlabel$};
      \node[above=5pt,interval label] at (\upper,\row) {$\rightlabel$};
    }
    \node[anchor=east,interval label] at (-2.6,-1.9) {$\bigcap_n I_n$};
    \draw[gray] (-1,-1.9) -- (1,-1.9);
    \node[interval closed] at (-1,-1.9) {};
    \node[interval closed] at (1,-1.9) {};
    \node[below=5pt,interval label] at (-1,-1.9) {$\ell_-=-1$};
    \node[below=5pt,interval label] at (1,-1.9) {$\ell_+=1$};
  \end{scope}
\end{tikzpicture}

  \caption[Tail bounds and extreme limit points.]{Terms and tail bounds
    for $x_n=(-1)^n(1+1/n)$. The first twelve terms are shown on the left.
    On the right, $I_n=[\alpha_n,\beta_n]$ contains the entire $n$th tail,
    including terms beyond those plotted. Its endpoints approach the two
    limit points $-1$ and $1$. The limiting interval bounds the limit
    points; it need not consist entirely of them.}
  \label{fig:tail-bounds-and-limit-points}
\end{figure*}

The tail infima and suprema need not be terms of the original sequence.
To show that their limits are limit points of $(x_n)$, we must select
actual terms approaching them. The following argument makes that
connection explicit for every bounded real sequence
\cite{Rodriguez2020Limsup}.
\begin{proposition}
  \label{thm:analysis-i:extreme-limit-points}
  The numbers $\ell_-$ and $\ell_+$ are respectively the smallest and
  largest limit points of a bounded real sequence $(x_n)$. Every limit
  point $\ell$ therefore satisfies $\ell_-\leq\ell\leq\ell_+$.
\end{proposition}
\begin{proof}
  Suppose $x_{n_k}\to\ell$. For each fixed $n$, all sufficiently late
  subsequence indices satisfy $n_k\geq n$. Thus
  $\alpha_n\leq x_{n_k}\leq\beta_n$. Taking limits first in $k$ and then
  in $n$ gives $\ell_-\leq\ell\leq\ell_+$.

  To reach the upper endpoint, set $n_0=0$. At step $k$, discard the
  terms through index $n_{k-1}$. The supremum of the remaining tail
  is $\beta_{n_{k-1}+1}$, so we can choose an actual term with
  $n_k>n_{k-1}$ and
  \[
    \beta_{n_{k-1}+1}-\frac1k<x_{n_k}\leq\beta_{n_{k-1}+1}.
  \]
  The indices increase, so both bounds tend to $\ell_+$. The squeeze
  theorem gives $x_{n_k}\to\ell_+$. Choosing increasing indices in the
  same way, with terms between the corresponding tail infimum and that
  infimum plus $1/k$, gives a subsequence converging to $\ell_-$.
  Hence both endpoints are limit points.
\end{proof}

For a bounded real sequence, showing that $L$ is its only limit point
is enough to establish $x_n\to L$. By
\autoref{thm:analysis-i:extreme-limit-points}, both $\ell_-$ and $\ell_+$
are limit points, so they must then equal $L$. Every original term
satisfies $\alpha_n\leq x_n\leq\beta_n$, and the squeeze theorem gives
$x_n\to L$. This is how information about subsequences controls the
whole sequence.

Equivalently, convergence occurs precisely when the two extreme limit
points agree.
If $x_n\to L$, then for every $\varepsilon>0$, all sufficiently late tails
lie in $[L-\varepsilon,L+\varepsilon]$. Their infima and suprema do too,
which proves that both limits equal $L$. The converse is the squeeze
argument just given.
This is the bounded-sequence version; more general upper and lower limits
will be studied with sequences later in the course.

The next exercise uses the elementary identities
\[
  \sin0=0,\qquad \sin\frac{2\pi}{3}=\frac{\sqrt3}{2},\qquad
  \sin\frac{4\pi}{3}=-\frac{\sqrt3}{2},
\]
together with $\sin(t+2\pi k)=\sin t$ for integers $k$.

% !TEX root = ../../hw03.tex
% Source: HW3 Intro to Math Analysis I (1).pdf, page 1, question 3.
\begin{exercise}
  \label{ex:hw03:oscillating-sequence}
  \solutionlink{hw03:oscillating-sequence}
  Consider the sequence
  \[
    a_n=\frac{1}{n}+\sin\left(\frac{2\pi n}{3}\right),
    \qquad n=1,2,\ldots.
  \]
  Find all limit points of $\{a_n\}$ and obtain
  \[
    \liminf_{n\to\infty}a_n,\quad\limsup_{n\to\infty}a_n.
  \]
\end{exercise}


\subsection{Which points can a set approach?}
\begin{definition}
  The \emph{closure} of $E\subseteq\R$ is
  \[
    \overline E=\{x\in\R\mid (x-\varepsilon,x+\varepsilon)\cap E
      \ne\varnothing\text{ for every }\varepsilon>0\}.
  \]
\end{definition}
A point of $E$ already belongs to its closure. A point $x$ is an
\emph{accumulation point} if every such neighborhood contains a point of
$E$ different from $x$. This extra condition distinguishes accumulation
points from general closure points.

\autoref{fig:closure-neighborhoods} compares these possibilities for
$E=(0,2)\cup\{3\}$. Every neighborhood of $0$ meets the interval, even
though $0\notin E$. A neighborhood of $3$ with radius less than $1$
meets $E$ only at $3$. At $5/2$, any radius less than $1/2$ misses $E$
entirely. The word \emph{every} in the definition is decisive.
\begin{figure*}[htbp]
  \centering
  % The ellipses are schematic set outlines; membership is read on the axis.
\begin{tikzpicture}[analysis diagram]
  \node at (6.8,2.1)
    {$E=(0,2)\cup\{3\},\qquad
      U_\varepsilon(x)=(x-\varepsilon,x+\varepsilon),\qquad
      0<\varepsilon_2<\varepsilon_1$};

  \foreach \shift/\testx/\smallradius/\panel in
    {0/0/0.4/1,5.3/3/0.3/2,10.6/2.5/0.3/3} {
    \begin{scope}[xshift=\shift cm]
      % The shaded lobe represents (0,2); the separate filled dot is {3}.
      \fill[gray!15] (1,0) ellipse[x radius=1cm,y radius=0.55cm];
      \begin{scope}
        \clip (1,0) ellipse[x radius=1cm,y radius=0.55cm];
        \fill[gray!40] (\testx,0)
          ellipse[x radius=\smallradius cm,y radius=0.4cm];
      \end{scope}
      \draw[gray] (1,0) ellipse[x radius=1cm,y radius=0.55cm];
      \node[analysis label] at (1.2,0.28) {$(0,2)$};

      % Two open neighborhoods with the same center and decreasing radius.
      \draw[dashed] (\testx,0)
        ellipse[x radius=0.8cm,y radius=0.85cm];
      \draw[densely dashed] (\testx,0)
        ellipse[x radius=\smallradius cm,y radius=0.4cm];
      \node[above=2pt,analysis label] at (\testx,0.85)
        {$\varepsilon_1$};
      \node[above=1pt,analysis label,fill=white] at (\testx,0.4)
        {$\varepsilon_2$};

      \draw[<->] (-0.95,0) -- (3.95,0);
      \draw[line width=1.2pt] (0,0) -- (2,0);
      \node[interval open] at (0,0) {};
      \node[interval open] at (2,0) {};
      \node[interval closed] at (3,0) {};
      \node[below=2pt,interval label] at (0,0) {$0$};
      \node[below=2pt,interval label] at (2,0) {$2$};
      \node[below=2pt,interval label] at (3,0) {$3$};

      \ifnum\panel=1
        \node at (1.5,1.5) {\textit{(a) Boundary point}, $x=0$};
        \node at (1.5,-1.18) {$0\in\overline E\setminus E$};
        \node[anchor=north,align=center,text width=4.5cm] at (1.5,-1.5)
          {Every $\varepsilon>0$ still reaches $E$.\\
           An accumulation point.};
      \else\ifnum\panel=2
        \node at (1.5,1.5) {\textit{(b) Isolated point}, $x=3$};
        \node at (1.5,-1.18) {$3\in E\subseteq\overline E$};
        \node[anchor=north,align=center,text width=4.5cm] at (1.5,-1.5)
          {Only $x$ remains when $\varepsilon<1$.\\
           Closure, but no accumulation.};
      \else
        \node at (1.5,1.5) {\textit{(c) Outside}, $x=\tfrac52$};
        \draw (2.5,-0.08) -- (2.5,0.08);
        \node[below=2pt,interval label] at (2.5,0) {$x$};
        \node at (1.5,-1.18) {$\tfrac52\notin\overline E$};
        \node[anchor=north,align=center,text width=4.5cm] at (1.5,-1.5)
          {A radius $\varepsilon<1/2$ misses $E$.\\
           One such radius excludes $x$.};
      \fi\fi
    \end{scope}
  }
\end{tikzpicture}

  \caption[Closure and shrinking neighborhoods.]{Schematic Venn views on
    a number line. The solid interval and filled point form $E$; its open
    endpoints are excluded. Dashed ellipses depict two shrinking
    $\varepsilon$-neighborhoods. The darker overlap in (a) persists as
    $\varepsilon$ decreases. Set membership is read on the axis; the
    ellipses add no points off the real line.}
  \label{fig:closure-neighborhoods}
\end{figure*}

\begin{proposition}
  A point $x$ belongs to $\overline E$ if and only if there is a sequence
  $(a_n)$ in $E$ converging to $x$.
\end{proposition}
\begin{proof}
  If $x\in\overline E$, choose $a_n\in E$ with $|a_n-x|<1/n$.
  The Archimedean property then gives $a_n\to x$. Conversely, if
  $a_n\to x$, every neighborhood of $x$ contains all sufficiently late
  $a_n$, and therefore meets $E$.
\end{proof}
For $a<b$, we have $\overline{(a,b)}=[a,b]$. Interior points already lie
in the interval, and every neighborhood of either endpoint meets it.
If $x<a$ or $x>b$, a neighborhood of radius less than its distance to the
nearest endpoint misses the interval. Thus no other point enters the closure.

\subsection{Rationals occur in every interval}
\begin{definition}
  A set $E\subseteq\R$ is \emph{dense in $\R$} if $\overline E=\R$.
  Equivalently, every nonempty open interval contains a point of $E$.
\end{definition}
For the equivalence, apply the neighborhood condition at the midpoint of
an open interval with radius half its length; in the other direction,
every neighborhood is itself an open interval.

\begin{theorem}
  \label{thm:analysis-i:rational-density}
  The rational numbers are dense in the real numbers.
\end{theorem}
\begin{proof}
  Given $x\in\R$ and $\varepsilon>0$, choose $N\in\N$ with
  $1/N<\varepsilon$. Set $m=\lfloor Nx\rfloor$. The definition of the
  floor gives $m\leq Nx<m+1$, and hence
  \[
    0\leq x-\frac{m}{N}<\frac{1}{N}<\varepsilon.
  \]
  Thus the rational number $m/N$ lies in the prescribed neighborhood
  of $x$. Every real number belongs to the closure of $\Q$.
\end{proof}
The grid in \autoref{fig:rational-grid} makes the estimate visible.
\begin{marginfigure}
  \centering
  \begin{tikzpicture}[interval diagram,x=13mm]
  \draw[<->] (-0.35,0) -- (2.7,0);
  \foreach \k in {0,1,2} {\draw (\k,-0.12) -- (\k,0.12);}
  \fill (0.65,0) circle[radius=1.6pt];
  \node[above=3pt,interval label] at (0.65,0) {$x$};
  \node[below,interval label] at (0,-0.15) {$\frac{m}{N}$};
  \node[below,interval label] at (1,-0.15) {$\frac{m+1}{N}$};
  \draw[<->] (0,-1.4) -- node[below] {$1/N$} (1,-1.4);
\end{tikzpicture}

  \caption[Approximating a real number on a rational grid.]{Schematic rational grid with spacing $1/N$. The floor chooses
    the left endpoint of the cell containing $x$.}
  \label{fig:rational-grid}
\end{marginfigure}
In particular, choosing $m_n=\lfloor nx\rfloor$ gives an explicit rational
sequence $m_n/n\to x$.

\subsection{Closure is relative to an ambient space}
Closure and completion answer related but different questions. Closure
adds points already present in a specified ambient space. Thus
$\overline{\Q}^{\,\R}=\R$, whereas $\overline{\Q}^{\,\Q}=\Q$.
A subset $E$ of an ambient space $X\subseteq\R$ is \emph{closed in $X$}
if it contains every point of its closure taken in $X$. Neighborhoods
in $X$ are intersections $(x-\varepsilon,x+\varepsilon)\cap X$.
For example, $[a,b]\cap\Q$ is closed in $\Q$. A rational point outside
it has a neighborhood missing $[a,b]$, while its real limit points
outside $\Q$ are not points of the ambient space.

A \emph{metric completion} instead embeds a metric space into a complete
one in which it is dense. For the usual distance $|x-y|$, $\R$ is a
completion of $\Q$.

More generally, $\overline E\subseteq\R$ is a completion of $E$ with its
usual distance. To check completeness, a Cauchy sequence $(y_n)$ in
$\overline E$ converges in $\R$ to some $y$. Given $\varepsilon>0$, choose
$n$ with $|y_n-y|<\varepsilon/2$, then choose $a\in E$ with
$|a-y_n|<\varepsilon/2$. The triangle inequality gives $|a-y|<\varepsilon$,
so $y\in\overline E$. Also, by the definition of closure, $E$ is dense
in $\overline E$. This explains why closing $(a,b)$ to $[a,b]$ also
completes that interval, while keeping the two notions distinct.

\subsection{Exercises on the missing limits in $\Q$}
The floor function and rational density let us build nested rational
intervals around a point missing from $\Q$. This complements the real
interval examples in \autoref{sec:analysis-i:completeness-equivalences}.

% !TEX root = ../../hw02.tex
% Homework 02, Exercise 3. Shared problem statement.
\begin{exercise}
  \label{ex:hw02:nested-rational-intervals}
  \solutionlink{hw02:nested-rational-intervals}
  Give an example of a sequence of nested nonempty bounded closed intervals
  $\{I_n\}$ in $\Q$ such that $I_{n+1}\subset I_n$ for each
  $n=1,2,\ldots$ but
  \[
    \bigcap_{n=1}^{\infty} I_n=\varnothing.
  \]
  Explain what goes ``wrong''.
\end{exercise}


\subsection{Geometric sums and factorial limits}
We need a simple estimate before the next two exercises.
Multiplying a finite geometric sum by $1-r$ and cancelling consecutive
terms gives
\[
  \sum_{k=0}^{N}r^k=\frac{1-r^{N+1}}{1-r}\qquad(r\ne1).
\]
For $0\leq r<1$, we have $r^N\to0$. If $r>0$, put $c=1/r-1>0$.
Induction gives $(1+c)^N\geq1+Nc$, and the Archimedean property shows
that $r^N\leq1/(1+Nc)\to0$. The case $r=0$ is immediate.
Taking limits in the finite identity gives
\begin{equation}
  \sum_{k=0}^{\infty}r^k=\frac1{1-r}\qquad(0\leq r<1).
  \label{eq:analysis-i:geometric-sum}
\end{equation}
Here, and in the series below, an infinite sum means the limit of its
finite partial sums. These particular sums need no general series theory.

\subsection{Completeness justifies positional notation}
We can now justify the representation facts used for cardinality in
\autoref{sec:analysis-i:power-sets}.
\begin{proof}[Proof of \autoref{thm:analysis-i:positional-representation}]
  For digits $a_k\in\{0,\ldots,b-1\}$, the partial sums
  $s_n=\sum_{k=1}^{n}a_k/b^k$ increase and satisfy
  $0\leq s_n\leq1-b^{-n}<1$ by the finite geometric identity.
  Monotone convergence supplies a limit in $[0,1]$.

  Conversely, let $0\leq x<1$, put $m_0=0$, and set
  $m_n=\lfloor b^nx\rfloor$. The floor inequalities give
  $bm_{n-1}\leq m_n\leq bm_{n-1}+b-1$.
  Thus $a_n=m_n-bm_{n-1}$ is an allowed digit, and telescoping gives
  \[
    \sum_{k=1}^{n}\frac{a_k}{b^k}=\frac{m_n}{b^n},\qquad
    0\leq x-\frac{m_n}{b^n}<b^{-n}\longrightarrow0.
  \]
  This constructs an expansion of $x$. Its digits cannot be eventually
  all $b-1$. If all digits after position $N$ were $b-1$, the geometric
  tail would give $x=(m_N+1)/b^N$, contrary to $b^Nx<m_N+1$.
  The endpoint $x=1$ is represented by taking every digit to be $b-1$.

  Finally, suppose two strings first differ at position $j$, with the
  first having the larger digit. Its advantage there is at least $b^{-j}$.
  The largest possible opposing tail is
  \[
    \sum_{k>j}\frac{b-1}{b^k}=b^{-j}.
  \]
  Equal values therefore require a digit difference of exactly one,
  followed by all zeros in the first string and all $b-1$ in the second.
  These are precisely the two forms of a terminating expansion.
  Excluding eventual strings of $b-1$ gives uniqueness in $[0,1)$.
\end{proof}

\subsection{The irrationality of $e$}
Set $s_n=\sum_{k=0}^{n}1/k!$. These partial sums increase.
For $k\geq2$, the bound $k!\geq2^{k-1}$ gives
\[
  s_n\leq2+\sum_{k=2}^{n}\frac1{2^{k-1}}<3.
\]
The monotone convergence theorem therefore supplies a real limit, denoted by
\[
  e=\sum_{k=0}^{\infty}\frac1{k!}.
\]
The alternative formula $e=\lim_{n\to\infty}(1+1/n)^n$ was also noted
in the lecture. Its equivalence with the series formula is deferred;
the series definition is enough for the present argument.

\begin{proposition}
  \label{thm:analysis-i:e-irrational}
  The number $e$ is irrational.\marginnote[-30pt]{What can we say about
    $\pi+e$ or $\pi/e$? Their irrationality is still unknown~\cite{WeissteinPi}.
    Knowing that $\pi$ and $e$ are individually transcendental does not
    settle either question.

    These questions matter because other proofs and theorems may require
    rationality or irrationality as a hypothesis. For example, irrationality
    of $\pi/e$ would imply that $a\pi+be=0$ with $a,b\in\Q$ forces
    $a=b=0$. Irrationality of $\pi+e$ would rule out $\pi+e=q$ for every
    $q\in\Q$. Such results exclude exact rational relations and supply
    hypotheses for later arguments. Until proved, any argument depending
    on these exclusions remains conditional.}
\end{proposition}
\begin{proof}
  Suppose $e=p/q$ for integers $p,q$ with $q>0$, and choose
  $n\geq\max\{q,2\}$. Both $n!e$ and $\sum_{k=0}^{n}n!/k!$ are integers.
  Their difference is
  \[
    R_n=n!\sum_{k=n+1}^{\infty}\frac1{k!}
       =\sum_{j=1}^{\infty}
         \frac1{(n+1)(n+2)\cdots(n+j)}.
  \]
  All terms are positive, and replacing each denominator factor by $n+1$
  gives, by \autoref{eq:analysis-i:geometric-sum},
  \[
    0<R_n<\sum_{j=1}^{\infty}\frac1{(n+1)^j}
          =\frac1n<1.
  \]
  The comparison is strict because the terms with $j\geq2$ are strictly
  smaller. Thus an integer lies strictly between zero and one, a
  contradiction.
\end{proof}
The factorial partial sums therefore give a bounded set of rationals
with an irrational real supremum. The next exercise asks for its bounds
and for the distinction between taking the supremum in $\R$ and in $\Q$.

% !TEX root = ../../hw02.tex
% Homework 02, Exercise 5. Shared problem statement.
\begingroup
\renewcommand{\labelenumi}{(\alph{enumi})}
\begin{exercise}
  \label{ex:hw02:factorial-sequence}
  \solutionlink{hw02:factorial-sequence}
  Consider the sequence
  \[
    x_n=\sum_{k=0}^{n}\frac{1}{k!},\quad n=1,2,\ldots,
  \]
  in $\Q$ and define $S=\{x_1,x_2,\ldots,x_n,\ldots\}$ which is a subset of $\Q$.
  \begin{enumerate}
    \item Show that $S$ is a bounded set.
    \item Find $a=\inf S$ and $b=\sup S$.
    \item Is it true that $a$ and $b$ both lie in $\Q$?
      Explain why and why not.
  \end{enumerate}
\end{exercise}
\endgroup


\subsection{Exercise on small interval covers}
The next question combines countable enumeration, rational density, and
\autoref{eq:analysis-i:geometric-sum} to construct interval covers
of small total length. Small covering length and metric incompleteness
describe different properties; completeness must be checked using Cauchy
sequences and their limits.

% !TEX root = ../../hw03.tex
% Source: HW3 Intro to Math Analysis I (1).pdf, page 2, question 5.
% The source's quoted terms "sparse" and "incompleteness" are retained.
% Small covering length alone is not a criterion for metric incompleteness.
\begingroup
\renewcommand{\labelenumi}{(\alph{enumi})}
\begin{exercise}
  \label{ex:hw03:rational-interval-cover}
  \solutionlink{hw03:rational-interval-cover}
  Let the rational numbers in $[0,1]$ be enumerated by
  $q_0,q_1,\ldots,q_n,\ldots$. Let $\varepsilon>0$ and cover this set
  $Q=\{q_n\}$ by
  \[
    U_\varepsilon=\bigcup_{n=0}^{\infty}I_n(\varepsilon),
  \]
  \[
    I_n(\varepsilon)=(q_n-\varepsilon 2^{-n},q_n+\varepsilon 2^{-n}),
    \qquad n=0,1,2,\ldots.
  \]
  That is, $Q\subset U_\varepsilon$.
  \begin{enumerate}
    \item Show that the sum of the lengths of the intervals
      $\{I_n(\varepsilon)\}$ satisfies the bound
      \[
        \sum_{n=0}^{\infty}|I_n(\varepsilon)|\leq4\varepsilon.
      \]
    \item Since $\varepsilon>0$ can be arbitrarily small, the set $Q$ can
      be covered by intervals of arbitrarily small total length indicating
      that $Q$ is a rather ``sparse'' subset of $[0,1]$. In view of this
      fact, give your explanation regarding the ``incompleteness'' of $Q$
      in $[0,1]$.
  \end{enumerate}
\end{exercise}
\endgroup

